% =============================================================================
%  Glossario e acronimi condivisi dai tre documenti.
%
%  I termini qui definiti sono gli stessi usati come nomi delle classi nel
%  codice: il glossario è il punto in cui il linguaggio del dominio viene
%  fissato una volta sola, per documenti e sorgenti insieme.
% =============================================================================

\chapter{Glossario}
\label{cap:glossario}

\section{Termini del dominio}

\begin{longtable}{@{}L{3.2cm} L{11.5cm}@{}}
\toprule
\textbf{Termine} & \textbf{Definizione} \\
\midrule
\endfirsthead
\toprule
\textbf{Termine} & \textbf{Definizione} \\
\midrule
\endhead
\bottomrule
\endlastfoot

Utente &
Persona registrata alla piattaforma. Consulta i tutorial, svolge i quiz,
lascia feedback e consegue badge. Nel codice: classe \id{Utente}. \\

Amministratore &
Utente con privilegi di gestione: crea e modifica tutorial, quiz e obiettivi,
modera i feedback e amministra gli account. È un ruolo dell'entità
\id{Utente}, non un'entità distinta. \\

Ruolo &
Attributo dell'utente che ne determina i permessi. Assume i valori
\id{utente} e \id{admin}. Nel codice: enumerazione \id{Ruolo}. \\

Tutorial &
Unità didattica composta da titolo, immagine di copertina, contenuto testuale
e categoria. È l'oggetto centrale del catalogo. \\

Categoria &
Classificazione tematica di un tutorial. L'insieme è chiuso: Internet,
Social Media, Tecnologia, Sicurezza. \\

Quiz &
Prova di verifica associata a un solo tutorial, composta da domande a scelta
multipla. La correzione è responsabilità esclusiva del server. \\

Domanda &
Quesito a scelta multipla appartenente a un quiz, con esattamente una
risposta corretta fra le opzioni proposte. \\

Risposta &
Singola opzione di una domanda. L'indicazione di correttezza non lascia mai
il server prima della consegna del quiz. \\

Svolgimento &
Registrazione del tentativo di un utente su un quiz: esito, numero di
risposte esatte e data. Un utente conserva un solo svolgimento per quiz. \\

Soglia di superamento &
Frazione di risposte corrette necessaria a superare un quiz, pari al 70\%. \\

Feedback &
Valutazione da 1 a 5 con commento, lasciata da un utente su un tutorial.
Un utente può lasciarne al più uno per tutorial. \\

Valutazione media &
Media dei feedback ricevuti da un tutorial, mantenuta dal database e
trattata come sola lettura dal codice applicativo. \\

Obiettivo &
Traguardo raggiungibile superando un numero prestabilito di quiz. Al
raggiungimento viene assegnato il badge associato. \\

Badge &
Riconoscimento visivo assegnato a un utente al conseguimento di un
obiettivo. \\

Conseguimento &
Associazione fra un utente e un obiettivo raggiunto, con la data di
assegnazione. \\

Catalogo &
Insieme dei tutorial pubblicati, consultabile anche senza autenticazione. \\

Sessione &
Periodo in cui un utente è riconosciuto dal sistema, materializzato da un
token firmato conservato in un cookie \id{httpOnly}. \\

\end{longtable}

\section{Termini tecnici}

\begin{longtable}{@{}L{3.2cm} L{11.5cm}@{}}
\toprule
\textbf{Termine} & \textbf{Definizione} \\
\midrule
\endfirsthead
\toprule
\textbf{Termine} & \textbf{Definizione} \\
\midrule
\endhead
\bottomrule
\endlastfoot

DAO &
\term{Data Access Object}: oggetto che incapsula l'accesso a una tabella o a
un aggregato del database, isolando il resto del sistema dal linguaggio SQL. \\

DTO &
\term{Data Transfer Object}: struttura dati priva di comportamento usata per
trasferire informazioni oltre il confine del sistema. In Tech4All è l'unico
punto in cui le entità vengono serializzate verso il client. \\

Servizio &
Oggetto che realizza le regole di dominio di un sottosistema, coordinando uno
o più DAO. Non conosce il protocollo HTTP. \\

Middleware &
Funzione interposta nella catena di elaborazione di una richiesta HTTP, che
può interromperla o arricchirla. \\

Aggregato &
Insieme di entità con un unico punto d'accesso e un ciclo di vita comune.
In Tech4All quiz, domande e risposte formano un solo aggregato. \\

Hash &
Trasformazione irreversibile di un valore. Le password sono conservate solo
come hash, mai in chiaro. \\

Sanificazione &
Rimozione dai contenuti HTML dei costrutti eseguibili, applicata prima della
persistenza. \\

Transazione &
Sequenza di operazioni sul database che viene applicata per intero o non
applicata affatto. \\

Mock &
Oggetto che sostituisce una dipendenza reale durante i test, esponendone
l'interfaccia con un comportamento deciso dal test stesso. In Tech4All sono
sostituiti da mock il pool di connessioni e i DAO ricevuti dai servizi: la
suite non richiede quindi alcun database. \\

Copertura del codice &
Frazione del codice sorgente eseguita durante i test, misurata su istruzioni,
diramazioni, funzioni e righe. \\

\end{longtable}

\section{Acronimi}

\begin{longtable}{@{}L{2.4cm} L{12.3cm}@{}}
\toprule
\textbf{Sigla} & \textbf{Significato} \\
\midrule
\endfirsthead
\toprule
\textbf{Sigla} & \textbf{Significato} \\
\midrule
\endhead
\bottomrule
\endlastfoot

API      & \term{Application Programming Interface} \\
CI       & \term{Continuous Integration} \\
CORS     & \term{Cross-Origin Resource Sharing} \\
CRUD     & \term{Create, Read, Update, Delete} \\
DBMS     & \term{Database Management System} \\
DTO      & \term{Data Transfer Object} \\
DAO      & \term{Data Access Object} \\
FURPS+   & \term{Functionality, Usability, Reliability, Performance, Supportability} e vincoli aggiuntivi \\
HTTP     & \term{HyperText Transfer Protocol} \\
JSON     & \term{JavaScript Object Notation} \\
JWT      & \term{JSON Web Token} \\
ODD      & \term{Object Design Document} \\
RAD      & \term{Requirements Analysis Document} \\
REST     & \term{Representational State Transfer} \\
RF       & Requisito Funzionale \\
RNF      & Requisito Non Funzionale \\
SDD      & \term{System Design Document} \\
SQL      & \term{Structured Query Language} \\
UC       & \term{Use Case}, caso d'uso \\
UML      & \term{Unified Modeling Language} \\
XSS      & \term{Cross-Site Scripting} \\

\end{longtable}
